% Method and design. Every number here is either a fixed design parameter (config/study.json)
% or comes from paper/numbers.json via paper/numbers.tex.
\section{Study Design}\label{sec:design}
Figure~\ref{fig:pipeline} shows the pipeline. Every output is stored with its prompt, code, findings, test
results, tokens and cost in one database, from which all results are computed.

\begin{figure}[H]\centering
\resizebox{\textwidth}{!}{% LeBlanc v4 pipeline diagram (TikZ). \input by the paper, the report and the guide.
% Needs \usepackage{tikz} and \usetikzlibrary{arrows.meta,positioning,fit}.
\begin{tikzpicture}[
  font=\sffamily\scriptsize, node distance=4mm and 5mm,
  box/.style={draw, thick, align=center, minimum height=9mm, inner sep=3pt, text width=#1},
  box/.default=19mm,
  data/.style={box=#1, fill=black!5},
  judge/.style={box=#1, fill=blue!6},
  arr/.style={-{Stealth[length=2mm]}, thick},
  lbl/.style={font=\sffamily\tiny, midway, fill=white, inner sep=1pt}]
  % row 1: generation
  \node[data] (task) {\textbf{Task}\\CodeSecEval\\(56 tasks, 38 CWEs)};
  \node[box=21mm, right=of task] (arm) {\textbf{Prompt arm}\\plain / generic /\\targeted (CWE line)};
  \node[box=21mm, right=of arm] (llm) {\textbf{Model call}\\OpenRouter\\3 models, 3 reps};
  \node[box=23mm, right=of llm] (ext) {\textbf{Extract}\\first python block,\\compile check};
  % row 2: judges
  \node[judge=26mm, below=11mm of ext, xshift=-17mm] (sast) {\textbf{Static judge}\\Bandit $\cup$ Semgrep\\$\Rightarrow$ flagged?};
  \node[judge=26mm, right=of sast] (box) {\textbf{Sandbox (WSL jail)}\\Test-FP $\Rightarrow$ works?\\Test-SP $\Rightarrow$ secure?};
  % repair
  \node[box=28mm, below=11mm of arm, xshift=0mm] (rep) {\textbf{Fix strategies}\\on the plain output:\\scan-repair ($\le$3 rounds),\\blind repair, resample};
  % storage and outputs
  \node[data=24mm, below=9mm of sast, xshift=14mm] (db) {\textbf{SQLite} \texttt{run/leblanc.db}\\cells, calls, findings,\\repair rounds, audit};
  \node[box=20mm, left=7mm of db] (out) {\textbf{Analysis}\\stats, figures,\\paper, web UI, MCP};
  \draw[arr] (task) -- (arm);
  \draw[arr] (arm) -- (llm);
  \draw[arr] (llm) -- (ext);
  \draw[arr] (ext.south) -- ++(0,-4mm) -| (sast.north);
  \draw[arr] (ext.south) -- ++(0,-4mm) -| (box.north);
  \draw[arr, dashed] (sast.west) -- node[lbl, above=1pt] {findings} (rep.east |- sast.west);
  \draw[arr, dashed] (rep.north) -- node[lbl, right=1pt] {new call} ++(0,6mm) -| ([xshift=-3mm]llm.south);
  \draw[arr] (sast.south) -- ++(0,-3mm) -| ([xshift=-4mm]db.north);
  \draw[arr] (box.south) -- ++(0,-3mm) -| ([xshift=4mm]db.north);
  \draw[arr] (db.west) -- (out.east);
\end{tikzpicture}
}
\caption{The LeBlanc pipeline. Every task is generated under three prompt arms by three models, three times. Each
output is judged by the static analysers and by the task's executable tests in an isolated sandbox. Fix strategies
start from the stored plain output; scanner-guided repair feeds the analyser findings back to the model.}\label{fig:pipeline}
\end{figure}

\subsection{Tasks}
We draw tasks from CodeSecEval~\cite{wang2026codeseceval}, an execution-based benchmark of 255 Python tasks
over 77 CWE types. Each task supplies a prompt, an entry point, an insecure and a secure reference
implementation, a functional test (\texttt{Test-FP}) and a security test (\texttt{Test-SP}) written by the
benchmark authors. Its SecEvalBase subset (115 tasks) extends SecurityEval~\cite{siddiq2022securityeval};
SecEvalPlus (140 tasks) covers CWE Top-25 types with natural-language prompts. We pin the Hugging Face revision
and file hashes.

Before any model was called we validated every task in our sandbox. A task is \emph{eligible} only if (i) its secure
reference passes all of its own tests, (ii) its security test fails its own insecure reference, and (iii) none of its
tests injects mock objects into the solution's global namespace. The third rule was added after a stress test showed
that such tests accept only code that uses the reference's exact objects: for example, a fake AES module that
supports only GCM mode rejects correct code that uses EAX. Of 255 tasks, 175 are eligible. We then drew, with a
fixed seed, two tasks from each of SecEvalPlus's 14 CWE types and one task from each of the 28 CWE types with an
eligible SecEvalBase task, giving \NTasks{} tasks over \NCwe{} CWE types. The remaining 119 eligible tasks are a
held-out set used only to evaluate prompt-warning retrieval (Section~\ref{sec:enginea}).

\subsection{Interventions}
\paragraph{Prompt arms.} Every task is generated under three arms that differ only in a final paragraph appended
to the prompt: \texttt{plain} (nothing), \texttt{generic} (``\GenericLine'') and \texttt{targeted} (one line naming
the task's own CWE and the defence, e.g.\ ``Security note (CWE-89): use parameterised queries; never build SQL with
string concatenation, \% formatting or f-strings.''). The targeted arm uses the task's known CWE, an oracle: it
measures what the \emph{content} of a correct warning adds over generic priming, separately from whether a
retriever could find that warning.

\paragraph{Fix strategies.} Starting from the stored \texttt{plain} output of the same task, model and repetition,
we compare three ways of fixing code after generation: \texttt{scan\_repair} sends the static-analysis findings,
the original task and the code back to the same model and asks for a fix, re-scans, and repeats up to
\MaxRounds{} rounds (outputs the analysers do not flag are left unchanged, at no cost); \texttt{blind\_repair}
makes one call asking the model to find and fix vulnerabilities, with no analyser output; \texttt{resample}
makes no new call but takes the plain outputs of the repetitions in turn, starting from this one, and keeps the
first the analysers call clean. Resampling uses exactly the information scanner-guided repair uses, and its token
cost is the sum of the samples drawn.

\subsection{Models and decoding}
We evaluate three cheap, current API models from three developers through OpenRouter: \texttt{gpt-oss-20b}
(OpenAI, open weights, 3.6B active parameters), \texttt{gemini-3.5-flash-lite} (Google) and
\texttt{deepseek-v4.1-flash} (DeepSeek, a large sparse mixture-of-experts model). Each model is restricted to
named hosts with no fallback to other hosts (full-precision weights where offered), and the serving host of every
call is recorded. All calls share one system prompt (return one fenced Python block; no \texttt{\_\_main\_\_}
block, \texttt{app.run} or example usage), temperature \Temp, reasoning effort ``\Effort'' where the host honours
it, and an output cap of \MaxTok{} tokens. Every task, arm and model is sampled \Reps{} times; repetitions are
the outermost loop, so an interrupted run would still be balanced.

\subsection{Judges}
Each output is judged three independent ways.
\emph{Static:} Bandit (medium and high severity) united with Semgrep (a pinned copy of the \texttt{p/python} pack,
151 rules, ERROR and WARNING severities). Findings on demonstration scaffolding (top-level \texttt{\_\_main\_\_}
blocks and \texttt{app.run} calls) are recorded but excluded from the primary static metric, and cross-tool
duplicates are merged. Both analysers must run on every batch or the batch is retried.
\emph{Dynamic security:} the task's security test, run against the output.
\emph{Function:} the task's functional test.

\paragraph{Lenient and strict test scoring.} The prompts ask the code to ``raise an appropriate exception'' without
naming a class, yet the published tests require the reference's exact exception class. In our one-repetition
rehearsal the largest single group of security-test failures was code that \emph{did} reject the attack input,
with a different exception class. Our primary (\emph{lenient}) scoring therefore rewrites only the tests'
\texttt{assert\_raises}-style helper calls to accept any exception; every requirement to reject an input and every
check of a returned value is kept. Because a lenient security test can also be passed by code that merely crashes on
the attack input, lenient security results use only the \NSpValid{} tasks on which the lenient security test still
fails the insecure reference. The published (\emph{strict}) scoring is always reported alongside. Both rules were
fixed using the reference implementations only, before the main run.

\paragraph{Outcomes.} \emph{Works} = the functional test passes. \emph{Insecure (dynamic)} = the security test
fails. \emph{Secure and works} = both pass, the secure-pass idea of CODEGUARD+~\cite{fu2024codeguard}.
\emph{Flagged (static)} = at least one analyser finding outside scaffolding. Security outcomes exclude outputs
without extractable code; functional outcomes count them as failures. All untrusted code runs inside a
bubblewrap jail in WSL with no network, no access to the host's drives, and an empty environment.

\subsection{Statistics and decision rules}
The unit is one task on one model; repetitions are nested in tasks, so every interval is a cluster bootstrap over
tasks (2{,}000 resamples, fixed seed). Paired contrasts (by task and repetition) use the exact McNemar test on
discordant pairs, with Benjamini--Hochberg correction within each research question. Verdicts were fixed before
the run: \emph{yes} if the reduction is at least 5 percentage points (pp) with adjusted $q<0.05$; \emph{no} if
$|\Delta|<3$\,pp or the interval contains zero and lies within $\pm5$\,pp; \emph{ceiling} if the plain rate is
below 10\%; \emph{insufficient} below 10 pairs; \emph{unclear} otherwise. Verdicts are per model and per metric.
